\bookchapter{Capitolo 7: Analisi completa delle funzioni avanzate}

\chapterintro{La tecnologia evolve di giorno in giorno e anche le funzioni del telefono continuano a progredire. Comunicazione satellitare, comunicazione senza rete, rilevamento incidenti, tecnologia SparkLink... queste funzioni che sembrano fantascientifiche sono ormai entrate nella nostra vita. Scopriamo insieme queste tecnologie all'avanguardia.}

\booksection{Comunicazione satellitare: chiedere aiuto anche senza segnale}

La comunicazione satellitare è una grande innovazione nel campo dei telefoni degli ultimi anni. Quando ti trovi in zone remote, in montagna o in una foresta e il telefono non ha assolutamente segnale, la comunicazione satellitare può essere l'ancora di salvezza.

\booksubsection{Principio di funzionamento:}

Il telefono si collega direttamente a un satellite di comunicazione in orbita per inviare messaggi di soccorso. Attualmente supporta principalmente l'invio di SMS e informazioni sulla posizione, ma non le chiamate vocali.

\booksubsection{Modelli supportati:}

\begin{itemize}
  \item \textbf{Apple}: iPhone 14 serie e successivi, supporta l'invio di SOS di emergenza via satellite. La serie iPhone 17 supporta ulteriormente l'invio di iMessage e SMS via satellite, con copertura estesa a più aree.
  \item \textbf{Huawei}: la serie Mate 50 inizia a supportare i messaggi satellitari Beidou (solo testo), la serie Mate 60 Pro supporta inoltre le chiamate satellitari Tiantong (possibilità di chiamare), la serie Mate 70 potenzia ulteriormente le capacità di comunicazione satellitare.
  \item \textbf{Xiaomi/OPPO/vivo}: dal 2024 i modelli di punta supportano progressivamente la funzione di messaggi satellitari Beidou.
  \item \textbf{Tendenza di sviluppo}: la comunicazione satellitare si sta diffondendo da funzione esclusiva dei modelli di punta di fascia alta verso il basso; in futuro anche i modelli di fascia media potrebbero supportarla. Il sistema satellitare cinese Beidou è quello con il supporto più ampio sui telefoni nazionali.
\end{itemize}

\booksubsection{Scenari d'uso:}

\begin{itemize}
  \item Escursionismo, incidenti in montagna.
  \item Incidenti in mare.
  \item Interruzione delle comunicazioni terrestri causata da disastri naturali.
\end{itemize}

\booksubsection{Modalità d'uso:}

\begin{enumerate}
  \item Assicurati di trovarti in zona aperta, con vista del cielo.
  \item Attiva la funzione di comunicazione satellitare del telefono (solitamente nell'interfaccia di chiamata di emergenza).
  \item Segui le indicazioni a schermo per orientare il telefono verso il satellite.
  \item Invia il messaggio di soccorso preimpostato e la posizione.
\end{enumerate}

\begin{tcbinfo}{Curiosità}
La comunicazione satellitare richiede un'antenna hardware speciale sul telefono, non tutti i telefoni la supportano. Inoltre, di solito richiede un abbonamento a pagamento, ma la funzione di SOS di emergenza è generalmente gratuita.

\end{tcbinfo}

\booksection{Comunicazione senza rete: un nuovo modo di comunicazione a breve distanza}

Per “comunicazione senza rete” si intende la tecnologia che permette a due telefoni di comunicare direttamente in assenza di rete mobile e Wi-Fi.

\booksubsection{Principio tecnico:}

Utilizza Bluetooth o frequenze radio dedicate per stabilire una connessione punto a punto a breve distanza (in genere da poche centinaia di metri a pochi chilometri).

\booksubsection{Scenari applicativi:}

\begin{itemize}
  \item \textbf{Attività all'aperto}: comunicazione tra membri di un gruppo durante trekking o campeggio.
  \item \textbf{Soccorso in caso di disastro}: quando le comunicazioni sono interrotte dopo terremoti o inondazioni.
  \item \textbf{Grandi eventi}: concerti, eventi sportivi con molta gente e segnale debole.
\end{itemize}

\booksubsection{Piattaforme supportate:}

\begin{itemize}
  \item \textbf{Apple}: "NameDrop" e "Trova amici offline" di iOS 17.
  \item \textbf{Huawei}: l'app MeeTime supporta le chiamate senza rete.
  \item \textbf{Xiaomi}: alcune funzioni del Xiaomi Share Center.
\end{itemize}

\booksection{Rilevamento incidenti: salva la vita nei momenti cruciali}

Il rilevamento incidenti è una funzione che utilizza i sensori del telefono per rilevare gravi incidenti stradali e chiamare automaticamente i numeri di emergenza.

\booksubsection{Principio di funzionamento:}

Il telefono utilizza accelerometro, giroscopio, barometro, microfono e altri sensori per valutare in modo integrato se si è verificato uno schianto grave. Rilevato l'incidente, il telefono:

\begin{enumerate}
  \item Mostra un'interfaccia con conto alla rovescia chiedendo se stai bene.
  \item Se alla fine del conto alla rovescia non rispondi, chiama automaticamente i numeri di emergenza.
  \item Invia la tua posizione ai contatti di emergenza.
\end{enumerate}

\booksubsection{Modelli supportati:}

\begin{itemize}
  \item \textbf{Apple}: iPhone 14 serie e successivi.
  \item \textbf{Android}: alcuni modelli di fascia alta lo supportano progressivamente.
\end{itemize}

\booksubsection{Avvertenze:}

\begin{itemize}
  \item La funzione deve essere attivata per poter essere utilizzata.
  \item Può attivarsi per errore (ad esempio sulle montagne russe o durante attività fisica intensa), basta annullare tempestivamente.
  \item Dipende da sensori e algoritmi, non è affidabile al 100\%, ma è una tutela in più.
\end{itemize}

\booksection{Tecnologia SparkLink: la connettività wireless di nuova generazione}

SparkLink è una nuova tecnologia di comunicazione wireless sviluppata sotto la guida di Huawei, con l'obiettivo di sostituire il Bluetooth e il Wi-Fi in alcuni scenari. Su indicatori chiave eguaglia o supera le soluzioni esistenti: la velocità di interfaccia può raggiungere 900Mbps (molto superiore ai 50Mbps del Bluetooth), latenza fino a 20 $\mu$ s (il Bluetooth circa 15ms), una singola rete supporta fino a 4096 dispositivi connessi contemporaneamente (il Bluetooth in genere meno di 8).

\booksubsection{Caratteristiche tecniche:}

\begin{itemize}
  \item \textbf{Latenza ultra-bassa}: latenza end-to-end di soli 20 microsecondi (0,02ms), due ordini di grandezza inferiore al Bluetooth, la sincronizzazione di giochi e musica è quasi impercettibile.
  \item \textbf{Alta velocità}: velocità di picco di un singolo collegamento fino a 900Mbps, decine di volte superiore all'audio Bluetooth, può trasmettere audio lossless e persino video ad alta risoluzione.
  \item \textbf{Basso consumo}: consumo in standby paragonabile al Bluetooth, ma con maggiore efficienza di trasmissione, il consumo energetico per unità di dati è inferiore.
  \item \textbf{Ampia copertura + alta concorrenza}: la distanza di connessione è superiore del 50\% rispetto al Bluetooth, una singola rete supporta oltre mille dispositivi connessi contemporaneamente.
\end{itemize}

\booksubsection{Scenari applicativi:}

\begin{itemize}
  \item \textbf{Cuffie wireless}: latenza inferiore, perfetta sincronizzazione audio-video nei giochi.
  \item \textbf{Orologi intelligenti}: connessione più stabile.
  \item \textbf{Casa intelligente}: più dispositivi connessi contemporaneamente.
  \item \textbf{Interconnessione automotive}: connessione telefono-sistema di bordo più fluida.
\end{itemize}

\booksubsection{Dispositivi supportati:}

Attualmente prevalentemente dispositivi dell'ecosistema Huawei, come la serie Huawei Mate 60, alcune cuffie e orologi Huawei. In futuro saranno supportati da più marchi.

\textbf{Prospettive tecnologiche}\par

La tecnologia SparkLink rappresenta la direzione futura della comunicazione wireless. Con l'adesione di più produttori, potremmo vedere più dispositivi che supportano SparkLink e godere di un'esperienza wireless migliore.

\begin{tcbpractice}{Verifica: il tuo telefono supporta queste funzioni}
\begin{enumerate}
  \item Cerca "satellite" nelle impostazioni per vedere se ci sono opzioni relative alla comunicazione satellitare.
  \item Cerca "rilevamento incidenti" o "SOS di emergenza" e verifica se sono attivi.
  \item Se hai un telefono Huawei, verifica se supporta il collegamento ai dispositivi SparkLink.
  \item Impara a usare queste funzioni: potrebbero tornare utili nei momenti cruciali.
\end{enumerate}

\end{tcbpractice}
